
%\needspace{4\baselineskip}
\color{uniblue}{\subsection{Логика отключения (ЛО)}}
\color{black}

\begin{enumerate}
%%%%%%%%%%%%%%%%%%%%%%%%%%%%%%%%%%%%%%%%%%%%%%%%%%%%%%%%%%%  ОБЩИЕ СВЕДЕНИЯ %%%%%%%%%%%%%%%%%%%%%%%%%%%%%%%%%%%%%%%%%%%%%%%%

\item Общие сведения

\begin{enumerate}

\item
Функция логики отключения выключателя предназначена для подачи аварийных команд отключения на выключатель при действии устройств РЗА. В устройстве логика представлена функциональным блоком «ЛО».

\item
Функциональный блок состоит из следующих функций:
\begin{itemize}
\item логика отключения от релейной защиты (ЛО РЗ); 
\item логика отключения ВНР (ЛО ВНР);
\item Логика запрета АВР (ЗАВР ЛО).
\end{itemize}

\item
Функциональный блок может быть полностью выведен из работы оперативно путем активации сигнала <<ОВ ЛО>>.
\item
В таблице \ref{lorz:tbl1} приведены параметры, необходимые для настройки ЛО.


\end{enumerate}

%%%%%%%%%%%%%%%%%%%%%%%%%%%%%%%%%%%%%%%%%%%%%%%%%%%%%%%%% Логика отключения от релейной защиты (ЛО РЗ) %%%%%%%%%%%%%%%%%%%%%%%%%%%%%%%%%%%%%%%%%%%%%%%%

\item Логика отключения от релейной защиты (ЛО РЗ)

\begin{enumerate}

\item
Функция <<ЛО РЗ>> формирует сигналы аварийного отключения (<<ЛО РЗ: Отключение>>, <<ЛО РЗ: Отключить>>) при срабатывании функций релейной защиты в составе устройства, а также от четырех назначаемых сигналов внешнего отключения (<<Откл. внешнее 1-4>>). В функциональном блоке «ЛО» логика отключения от релейной защиты представлена функцией «ЛО РЗ».

\item
Сигнал <<ЛО РЗ: Отключение>> является суммарным сигналом отключения от отдельных функций и существует пока активен хотя бы один сигнал срабатывания защитных функций.

\item
Сигнал <<ЛО РЗ: Отключить>> воздействует на выходные реле для отключения выключателя, длительность сигнала задается уставкой Т1.

\item
Ввод функции в работу на этапе параметрирования устройства осуществляется программным переключателем SGF1 <<Ввод\_функции>>. Информация о введенной функции отображается выходным сигналом <<ЛО РЗ: Ввод>>.

\item
Функция может быть выведена из работы оперативно путем активации сигнала <<ОВ ЛО РЗ>>, что характеризуется активным состоянием сигнала <<ЛО РЗ: Оперативный вывод>> на выходе функции.

\item
Алгоритм работы функции выполнен в соответствии с рисунком \ref{lorz:img1}.

\vspace{3mm}
\begin{figure}[H]
\centering
\includegraphics[width=0.85\textwidth,height=0.85\textheight,keepaspectratio]{\fbpath/412. ЛО СВ/_latex/img/lorz.pdf}
\captionof{figure}{Функциональная схема ЛО РЗ}\label{lorz:img1}
\end{figure}

\end{enumerate}

%%%%%%%%%%%%%%%%%%%%%%%%%%%%%%%%%%%%%%%%%%%%%%%%%%%%%%%%%  Логика отключения ВНР (ЛО ВНР)  %%%%%%%%%%%%%%%%%%%%%%%%%%%%%%%%%%%%%%%%%%%%%%%%

\item  Логика отключения ВНР (ЛО ВНР) 

\begin{enumerate}

\item\sloppy
ЛО ВНР формирует сигналы аварийного отключения (<<ЛО ВНР: Отключение>>, <<ЛО ВНР: Отключить>>) при срабатывании внешней функции ВНР <<Отключение от ВНР>>. В функциональном блоке «ЛО» логика отключения ВНР представлена функцией «ЛО ВНР». 

\item
Сигнал <<ЛО ВНР: Отключить>> воздействует на выходные реле для отключения выключателя, длительность сигнала задается уставкой Т1.

\item
Ввод функции в работу на этапе параметрирования устройства осуществляется программным переключателем SGF1 <<Ввод\_функции>>. Информация о введенной функции отображается выходным сигналом <<ЛО ВНР: Ввод>>.

\item
Функция может быть выведена из работы оперативно путем активации сигнала <<ОВ ВНР>>, что характеризуется активным состоянием сигнала <<ЛО ВНР: Оперативный вывод>> на выходе функции.

\item
Алгоритм работы функции выполнен в соответствии с рисунком \ref{lovnr:img1}.

\vspace{3mm}
\begin{figure}[H]
\centering

\includegraphics[width=0.85\textwidth,height=0.85\textheight,keepaspectratio]{\fbpath/412. ЛО СВ/_latex/img/lovnr.pdf}
\captionof{figure}{Функциональная схема ЛО ВНР}\label{lovnr:img1}
\end{figure}

\end{enumerate}
%%%%%%%%%%%%%%%%%%%%%%%%%%%%%%%%%%%%%%%%%%%%%%%%%%%%%%%%% КОНЕЦ Логика отключения ВНР (ЛО ВНР) %%%%%%%%%%%%%%%%%%%%%%%%%%%%%%%%%%%%%%%%%%%%%%%%


%%%%%%%%%%%%%%%%%%%%%%%%%%%%%%%%%%%%%%%%%%%%%%%%%%%%%%%%% Запрет АВР %%%%%%%%%%%%%%%%%%%%%%%%%%%%%%%%%%%%%%%%%%%%%%%%

\item Логика запрета АВР (ЗАВР ЛО)

\begin{enumerate}

\item
Запрет АВР может быть необходим в определённых ситуациях для предотвращения повторных срабатываний и возможных повреждений оборудования. В функциональном блоке «ЛО» логика запрета АВР представлена функцией «ЗАВР ЛО».

\item
Сигнал <<ЗАВР: Запрет АВР>> формируется, если присутствует сигнал срабатывания защит на любом из входов, которые заведены под запрет АВР с помощью выбора уставки программного переключателя. Контроль всех входных сигналов внешнего отключения выполняется соответствующими уставками (SGF1 \ldots SGF4).

\item
Ввод функции в работу осуществляется при введенной функции АВР в работу (<<АВР: Ввод>>). Информация о введенной функции отображается выходным сигналом <<ЛОВ / ЗАВР: Ввод>>.

\item
Функция  может быть выведена из работы оперативно путем активации сигнала <<ОВ ЛОВ>>, что характеризуется активным состоянием сигнала <<ЛОВ / ЗАВР: Оперативный вывод>> на выходе функции.

\item
Алгоритм работы функции выполнен в соответствии с рисунком \ref{zavr:img1}. 

\vspace{3mm}
\begin{figure}[H]
\centering

\includegraphics[width=0.8\textwidth,height=0.8\textheight,keepaspectratio]{\fbpath/412. ЛО СВ/_latex/img/zavr.pdf}
\captionof{figure}{Функциональная схема логики ЗАВР}\label{zavr:img1}
\end{figure}

\end{enumerate}
%%%%%%%%%%%%%%%%%%%%%%%%%%%%%%%%%%%%%%%%%%%%%%%%%%%%%%%%% КОНЕЦ Запрет АПВ %%%%%%%%%%%%%%%%%%%%%%%%%%%%%%%%%%%%%%%%%%%%%%%%


\end{enumerate}
